\documentclass[11pt]{article}

\usepackage[a4paper,margin=2.7cm]{geometry}
\usepackage{amsmath,amssymb,amsthm,mathtools}
\usepackage{enumitem}
\usepackage{hyperref}

\newtheorem{theorem}{Theorem}[section]
\newtheorem{proposition}[theorem]{Proposition}
\newtheorem{lemma}[theorem]{Lemma}
\newtheorem{remark}[theorem]{Remark}

\newcommand{\C}{\mathbb C}
\newcommand{\R}{\mathbb R}
\newcommand{\Z}{\mathbb Z}
\newcommand{\dd}{\,d}
\newcommand{\pbar}{\partial_{\bar z}}
\newcommand{\Log}{\operatorname{Log}}

\title{A Parametric Asymptotic Route for the One-Boundary Hermite Rigidity Problem}
\author{}
\date{}

\begin{document}
\maketitle

\section{Boundary moments}

This note isolates the boundary-parametrization route for the simply connected
Hermite rigidity theorem.  It deliberately does not use the holomorphic auxiliary
function argument from Section~4 of \texttt{main.tex}.  The input is only the
moment identity coming from the Hermite eigenfunction equation and the standard
free-boundary consequence that the one boundary component is real analytic.

Let \(U\subset\C\) be bounded and simply connected, with real-analytic boundary
\(\Gamma=\partial U\), positively oriented.  Assume the Hermite eigenfunction
identity gives
\begin{equation}\label{eq:bulk-moments}
  \int_U z^n |z|^{2k}e^{-\pi |z|^2}\dd A(z)=0,
  \qquad n\ge1 .
\end{equation}
Set
\begin{equation}\label{eq:Phi-def}
  \Phi_k(s):=\int_0^s t^k e^{-\pi t}\dd t .
\end{equation}
Then
\[
  \pbar\!\left(z^{n-1}\Phi_k(|z|^2)\right)
  =
  z^n|z|^{2k}e^{-\pi|z|^2},
\]
and Cauchy--Green gives
\begin{equation}\label{eq:boundary-moments}
  M_n:=\int_{\Gamma}z^{n-1}\Phi_k(|z|^2)\dd z=0,
  \qquad n\ge1 .
\end{equation}

Let
\[
  \gamma:\R/2\pi\Z\to\Gamma
\]
be a real-analytic parametrization.  It extends holomorphically to a strip
\[
  \mathcal S_\epsilon=\{t\in\C/2\pi\Z:|\Im t|<\epsilon\}.
\]
Define the reflected analytic continuation
\[
  \gamma^*(t):=\overline{\gamma(\bar t)}.
\]
On the real axis, \(\gamma^*(t)=\overline{\gamma(t)}\).  Thus
\begin{equation}\label{eq:param-moment}
  M_n=\int_0^{2\pi}
  \gamma(t)^{n-1}
  E(t)\gamma'(t)\dd t,
  \qquad
  E(t):=\Phi_k(\gamma(t)\gamma^*(t)).
\end{equation}

The useful normalization is obtained by integrating by parts:
\begin{equation}\label{eq:ibp}
  M_n
  =
  -\frac1n\int_0^{2\pi}\gamma(t)^n E'(t)\dd t .
\end{equation}
This is the right starting point for asymptotics.  The older real-maximum
heuristic misses this cancellation: the amplitude
\(E(t)\gamma'(t)/\gamma(t)\) has a built-in derivative structure.

\begin{lemma}\label{lem:constant-radius}
If \(E'\equiv0\) on the real axis, then \(\Gamma\) is a circle centered at the
origin.
\end{lemma}

\begin{proof}
Since \(\Phi_k'(s)=s^k e^{-\pi s}>0\) for \(s>0\), the identity
\[
  E'(t)=\Phi_k'(|\gamma(t)|^2)\frac{\dd}{\dd t}|\gamma(t)|^2
\]
on the real axis implies \(\frac{\dd}{\dd t}|\gamma(t)|^2=0\).  Hence
\(|\gamma|\) is constant on \(\Gamma\), so \(\Gamma\) is a circle centered at
the origin.
\end{proof}

\section{Explicit saddle asymptotics}

Assume from now on that \(\Gamma\) is not a centered circle.  Then
\(E'\not\equiv0\).  Choose a branch
\[
  \Psi(t):=\Log \gamma(t)
\]
on the relevant lifted strip.  Equation~\eqref{eq:ibp} becomes
\begin{equation}\label{eq:saddle-integral}
  M_n=-\frac1n\int e^{n\Psi(t)}E'(t)\dd t .
\end{equation}

\begin{proposition}[Dominant saddle contribution]\label{prop:saddle}
Suppose the real cycle can be deformed, in the analytic continuation domain of
\(\gamma,\gamma^*,E\), to steepest-descent arcs through finitely many dominant
saddles \(\tau_1,\dots,\tau_L\).  At each \(\tau_j\), assume
\[
  \Psi(t)=\Psi(\tau_j)+\lambda_j(t-\tau_j)^{p_j}
  +O((t-\tau_j)^{p_j+1}),
  \qquad p_j\ge2,\quad \lambda_j\ne0,
\]
and
\[
  E'(t)=b_j(t-\tau_j)^{r_j}+O((t-\tau_j)^{r_j+1}),
  \qquad b_j\ne0,\quad r_j\ge0.
\]
Assume also that the corresponding exponential scales dominate all remaining
parts of the deformed contour.  Let
\[
  \rho:=\max_j |\gamma(\tau_j)|.
\]
Among the saddles with \(|\gamma(\tau_j)|=\rho\), keep only those for which the
power
\[
  \alpha_j:=1+\frac{r_j+1}{p_j}
\]
is minimal; call this set \(J_*\).  Then
\begin{equation}\label{eq:saddle-asymptotic}
  M_n
  =
  \rho^n n^{-\alpha_*}
  \left(
    \sum_{j\in J_*} C_j e^{in\theta_j}
  \right)
  +o(\rho^n n^{-\alpha_*}),
\end{equation}
where
\[
  \gamma(\tau_j)=\rho e^{i\theta_j},
  \qquad
  \alpha_*=\min_{|\gamma(\tau_j)|=\rho}\alpha_j .
\]
The constants are explicitly
\begin{equation}\label{eq:Cj}
  C_j
  =
  -b_j\,\sigma_j\,
  \Gamma\!\left(\frac{r_j+1}{p_j}\right)
  \lambda_j^{-(r_j+1)/p_j},
\end{equation}
where \(\sigma_j\ne0\) is the difference of the two steepest-ray phase factors
chosen by the deformed oriented contour.
\end{proposition}

\begin{proof}
Near \(\tau_j\), put \(s=t-\tau_j\).  The local contribution to
\eqref{eq:saddle-integral} is
\[
  -\frac1n e^{n\Psi(\tau_j)}
  \int e^{n\lambda_j s^{p_j}}
  \bigl(b_js^{r_j}+O(s^{r_j+1})\bigr)\dd s .
\]
On a steepest ray, set \(u=n^{1/p_j}\lambda_j^{1/p_j}s\), with the branch fixed
by that ray.  This gives
\[
  -e^{n\Psi(\tau_j)}
  n^{-1-(r_j+1)/p_j}
  b_j\lambda_j^{-(r_j+1)/p_j}
  \int_{\mathcal R_j} e^{u^{p_j}}u^{r_j}\dd u
\]
plus lower powers of \(n^{-1/p_j}\).  The steepest-ray integral is the nonzero
phase factor
\[
  \sigma_j\,\Gamma\!\left(\frac{r_j+1}{p_j}\right),
\]
unless the oriented contour misses the local saddle contribution, which is
excluded by the dominance/accessibility hypothesis.  Summing the dominant
saddles and absorbing lower-order saddles and the deformed remainder gives
\eqref{eq:saddle-asymptotic}.
\end{proof}

\begin{lemma}\label{lem:exp-sum}
Let \(\lambda_1,\dots,\lambda_N\) be distinct complex numbers of modulus one.
If
\[
  \sum_{j=1}^N c_j\lambda_j^n\to0,
\]
then \(c_1=\cdots=c_N=0\).
\end{lemma}

\begin{proof}
Multiply by \(\overline{\lambda_m}^{\,n}\), average from \(1\) to \(N\), and let
\(N\to\infty\).  The Cesaro averages of
\((\lambda_j\overline{\lambda_m})^n\) vanish for \(j\ne m\), while the averaged
left-hand side still tends to zero.
\end{proof}

\begin{theorem}[Conditional asymptotic closure]\label{thm:conditional-saddle}
Assume the saddle hypotheses of Proposition~\ref{prop:saddle}, and assume that
the dominant phases in \(J_*\) have been grouped so that the resulting finite
exponential sum in \eqref{eq:saddle-asymptotic} is not identically zero.  Then
the moment identities \eqref{eq:boundary-moments} force \(\Gamma\) to be a
circle centered at the origin.
\end{theorem}

\begin{proof}
If \(\Gamma\) is not a centered circle, then \(E'\not\equiv0\) by
Lemma~\ref{lem:constant-radius}.  Proposition~\ref{prop:saddle} gives a nonzero
asymptotic for \(M_n\):
\[
  \rho^{-n}n^{\alpha_*}M_n
  \to
  \sum_{j\in J_*} C_j e^{in\theta_j}
  \quad\text{up to }o(1).
\]
Since \(M_n=0\) for every \(n\), the finite exponential sum must tend to zero.
By Lemma~\ref{lem:exp-sum}, all grouped coefficients vanish, contradicting the
nonzero dominant-sum hypothesis.  Hence \(\Gamma\) must be a centered circle.
\end{proof}

\section{Darboux singularity version}

The same computation can be made in the \(z\)-plane.  Let \(S\) denote the local
Schwarz function of the analytic boundary, and define in a collar of \(\Gamma\)
\[
  Q(z):=\Phi_k(zS(z)).
\]
Then \(Q(z)=\Phi_k(|z|^2)\) on \(\Gamma\), and
\[
  M_n=\int_\Gamma z^{n-1}Q(z)\dd z .
\]

\begin{proposition}[Pole contribution]\label{prop:pole}
Suppose the contour \(\Gamma\) can be pushed into \(U\) until it crosses finitely
many dominant pole singularities \(\alpha\) of \(Q\), and suppose the remaining
contour has exponentially smaller \(|z|\).  If near such a pole
\[
  Q(z)=\sum_{\ell=1}^{q_\alpha}
  a_{\alpha,-\ell}(z-\alpha)^{-\ell}+O(1),
\]
then its contribution to \(M_n\) is exactly
\begin{equation}\label{eq:pole-contribution}
  2\pi i
  \sum_{\ell=1}^{q_\alpha}
  a_{\alpha,-\ell}
  \binom{n-1}{\ell-1}
  \alpha^{\,n-\ell}.
\end{equation}
Thus the dominant pole set gives
\[
  M_n
  =
  \rho^n n^{q_*-1}
  \left(\sum_{|\alpha|=\rho}D_\alpha e^{in\arg\alpha}\right)
  +o(\rho^n n^{q_*-1}),
\]
after retaining the maximal pole order \(q_*\) among poles with
\(|\alpha|=\rho\).
\end{proposition}

\begin{proof}
This is the residue theorem applied to \(z^{n-1}Q(z)\).  For a pole term
\((z-\alpha)^{-\ell}\), the residue is
\[
  \frac{1}{(\ell-1)!}
  \frac{\dd^{\ell-1}}{\dd z^{\ell-1}}z^{n-1}\bigg|_{z=\alpha}
  =
  \binom{n-1}{\ell-1}\alpha^{n-\ell}.
\]
The stated asymptotic follows by keeping the largest \(|\alpha|\) and the
largest pole order at that modulus.
\end{proof}

\begin{remark}[What remains to make this unconditional]
The computation above is fully explicit once the dominant complex obstruction is
known: either accessible saddles of \(\Log\gamma\), or Darboux singularities of
the continued weighted Schwarz primitive \(Q\).  The missing theorem is not an
asymptotic calculation; it is a continuation theorem saying that every noncircular
free boundary arising from the Hermite moment problem has such a finite dominant
obstruction with nonzero grouped leading coefficient.  Real analyticity alone
does not provide that statement.  It gives a complex collar, but not a finite
singularity structure or a global meromorphic Schwarz function.
\end{remark}

\end{document}
